\documentclass{article}
 
\usepackage[ruled,vlined,linesnumbered]{algorithm2e}
\usepackage{geometry}
\usepackage{amsmath}
\DeclareMathOperator*{\argmax}{arg\,max}
\DeclareMathOperator*{\argmin}{arg\,min}
\geometry{a4paper, portrait, margin=2in}
\newenvironment{breakablealgorithm}
\SetAlFnt{\large}
\SetAlCapFnt{\large}
\SetAlCapNameFnt{\large}
\SetAlgoLined
\begin{document}

 
% The Monte Carlo Tree Search Algorithm



%%%%%%%%%%%%% THE FIRST ALGORITHM %%%%%%%%%%%%%%%%%%%%%%%%%%%%
\begin{algorithm}[H]


\SetKwInOut{Input}{Input}\SetKwInOut{Output}{Output}
\Input{The state, $s_0$, of the current root node, $n_0$ (s($n_0$))}
\Output{The action, $a$, that leads to the most optimum child node that has a desirable state}
\BlankLine
\SetKwProg{Fn}{Function}{}{}
\Fn{MCTS($s_0$)}{
create root node $n_0$ with state $s_0$\; 
\BlankLine

\While{within computational budget}{    
    
    \BlankLine
    \tcc{The $SELECTEXPAND$ function is used to select a child node according to a predefined utility function: ``UCT lean-LGRF'' in the case of this study. This selected child node $n_i$ is then added to the search tree in the process of expansion. The variated expansion stage goes a step further to add all the 3 possible children of the child node $n_i$.}
    $n_i \leftarrow SELECTEXPAND(s(n_0))$    \;
    \BlankLine
    \tcc{The $SIMULATEBACKPROPAGATE$ function is used to randomly simulate possible scenarios from the newly expanded node $n_i$. A reward value is calculated by executing each of the possible actions available from the newly expanded node. This is then stored in $\Delta$ and back-propagated up the tree to update the node values.}
    $\Delta \leftarrow SIMULATEBACKPROPAGATE(s(n_i))$    \;
}
\BlankLine

\tcc{The result of the overall search, $a(SELECT(n_0))$, is the action, $a$, that leads to the best child of the root node, $n_0$, where the exact definition of ``best'' is defined by the implementation. In the case of this study, ``best'' is the node with the highest discounted transaction throughput (for OLTP) or the lowest discounted response-time latency (for OLAP).}
\Return $a(SELECT(n_0))$ \;
}

\caption{The Basic Monte Carlo Tree Search Algorithm}
\end{algorithm}

%%%%%%%%%%%%%%%%%%%%%%%%%%%SECOND ALGORITHM%%%%%%%%%%%%%%%
\begin{algorithm}[H]


\SetKwInOut{Input}{Input}\SetKwInOut{Output}{Output}
\Input{The state, $s_0$, of the current root node, $n_0$ (s($n_0$))}
\Output{The child node $n_i$ that has been selected through either exploitation or exploration. In addition to $n_i$, all its 3 children are also specified as output.}
\BlankLine

\SetKwProg{Fn}{Function}{}{}
\Fn{SELECTEXPAND($n_0$)}{
\While{$n_0$ is non-terminal}{    
  \uIf{$n_0$ has been fully expanded}{
    \Return $\Delta \leftarrow SELECT(n_0)$    \;
  }
  \Else{
    \Return $EXPAND(n_0)$   \;
  }
}
}

\BlankLine
\BlankLine
\BlankLine

\SetKwProg{Fn}{Function}{}{}
\Fn{SELECT($n_0$)}{

\tcc{Determines the best child, $\Delta$, to select out of all the other possible children of node $n_0$. This is based on the value of the UCT lean-LGRF variation to the selection stage.}

\Return \( \displaystyle \argmax_{i \in children Of The Root Node} \Bigg(\frac{\gamma (n_i)}{x(n_i)} + lean\_LGRF(n_0, n_i) + \sqrt {c \times \frac{ \log (t)}{x(n_i)}}\Bigg) \) \;

}

\BlankLine
\BlankLine
\BlankLine

\SetKwProg{Fn}{Function}{}{}
\Fn{EXPAND($n_0$)}{
    perform $a \in untried$ $actions$ from $A(s(n_0))$ \;
    add new child $n_i$ to $n_{i-1}$
    \linebreak
    with $a(n_i) \leftarrow a$
    \linebreak
    and $s(n_i) \leftarrow f(s(n_{i-1}), a)$ \;
    add all children of $n_i$ to $n_i$ \;
    \Return $n_i$, $children\ of\ n_i$ \;
}

\caption{Variated Selection and Variated Expansion Stages of the Monte Carlo Tree Search Algorithm}
\end{algorithm}

\pagebreak
%%%%%%%%%%%%%%%%%%%%%%%%%%%%THIRD ALGORITHM%%%%%%%%%%%%%%%
\begin{algorithm}[H]


\SetKwInOut{Input}{Input}\SetKwInOut{Output}{Output}
\Input{The state, $s_i$, of the current node, $n_i$ (s($n_i$))}
\Output{The reward value of node $n_i$}

\SetKwProg{Fn}{Function}{}{}
\Fn{SIMULATEBACKPROPAGATE($s(n_i)$)}{
    \While{$A(s(n_i))$ is not empty and $x(n_i)<$1,000,000}{
        perform $a \in A(s(n_i))$ at random \;
        $s(n_i) \leftarrow f(s(n_{i-1}),a)$ \;
        \tcc{$x(n_i)$ is the number of times node $n_i$ has been traversed}
        $x(n_i) \leftarrow x(n_i) + 1$ \;
        \uIf{$result$ is a win}{
        \tcc{The number of simulated wins that have occurred when node $n_i$ was traversed; where a ``win'' is defined by the implementation. In the case of this study, ``win'' is the node with the highest discounted transaction throughput (for OLTP) or the least response-time latency (for OLAP).}
            $w_i \leftarrow w_i$ + 1 \;
        }
    }
        SIMULATEBACKPROPAGATE(for all the children of $n_i$) \;
        BACKPROPAGATE($n_i$) \;
}

\BlankLine
\BlankLine
\BlankLine

\SetKwProg{Fn}{Function}{}{}
\Fn{BACKPROPAGATE($n_i$)}{
    \While{$n_i$ is not $null$}{
        update $rewardValueOfNode$ $n_i$ \;
        \tcc{$w_i$ = discounted transaction throughput (for OLTP) or the least response-time latency (for OLAP) for node $n_i$ and all of its children}
        update $w_i$ \;
        \tcc{$x_i$ = number of traversals (visits) made on node $n_i$ and all of its children}
        update $x_i$ \;
        update $rewardValueOfNode$ $n_i = \dfrac{\gamma (n_i)}{x(n_i)}$ \;
        $n_i \leftarrow$ parent of $n_{i-1}$ \;
    }
}
\caption{Variated Simulation and Back-Propagation Stages of the Monte Carlo Tree Search Algorithm}
\end{algorithm}

\end{document}